\documentclass[11pt]{article}

\usepackage[margin=1in]{geometry}
\usepackage{graphicx}
\usepackage{booktabs}
\usepackage{array}
\usepackage{xcolor}
\usepackage{hyperref}
\usepackage[numbers,sort&compress]{natbib}
\usepackage{caption}
\usepackage{float}

\graphicspath{{../../results/manuscript_assets/}{../../results/figures/esm1v_ensemble/}}

\hypersetup{
  colorlinks=true,
  linkcolor=blue,
  citecolor=blue,
  urlcolor=blue
}

\title{Zero-shot prioritization and mechanistic annotation of MAPT/Tau missense variants}
\author{MAPT Zero-Shot Project}
\date{Draft reading version}

\begin{document}
\maketitle

\begin{abstract}
MAPT/Tau missense variants are difficult to interpret because many possible
substitutions lack clinical or experimental labels. We generated a complete
zero-shot atlas for the 441-aa Tau-F / hTau40 isoform by enumerating all 8379
single amino-acid substitutions and scoring them with an ESM-1v ensemble. We
added Tau-specific mechanistic annotations and integrated ClinVar and
AlphaMissense using strict reference-residue coordinate quality control. The
atlas highlights higher zero-shot sensitivity in Tau microtubule repeat regions
and prioritizes ClinVar variants of uncertain significance for follow-up. Because
the strict ClinVar binary benchmark is small, this work should be interpreted as
a reproducible prioritization atlas, not as a standalone clinical diagnostic
model.
\end{abstract}

\noindent\textbf{Keywords:} MAPT; Tau; missense variants; zero-shot prediction;
protein language models; ESM-1v; AlphaMissense; ClinVar.

\section{Introduction}

Protein-coding missense variants are common in human genetic data, but many are
difficult to interpret. A missense variant changes one amino acid in a protein.
Some such changes are harmless, while others disrupt protein function, alter
protein interactions, or increase disease risk. For many disease-associated
genes, the number of possible missense variants is much larger than the number of
variants that have been experimentally or clinically characterized.

MAPT encodes the microtubule-associated protein Tau, a neuronal protein involved
in microtubule binding and stabilization. Tau is also central to tauopathies,
where abnormal Tau pathology is a defining feature. In inherited frontotemporal
dementia and parkinsonism linked to chromosome 17, missense and splice-site
mutations in MAPT provided direct genetic evidence that Tau dysfunction can cause
neurodegeneration \citep{hutton1998tau,spillantini2000tau}.

The biological interpretation of MAPT variation is complicated by Tau isoforms.
Adult human brain expresses multiple Tau isoforms generated by alternative
splicing. These isoforms differ in N-terminal inserts and in the number of
microtubule-binding repeats. The 441-amino-acid 2N4R Tau-F / hTau40 isoform is a
commonly used reference for Tau biology, but external databases may report MAPT
variants using different transcript or protein coordinate systems. This makes
strict coordinate checking essential before combining clinical labels, external
predictor scores, and protein-model outputs \citep{spillantini2000tau}.

ClinVar provides clinical interpretations of human variants and is useful for
benchmarking computational predictions \citep{landrum2018clinvar}. However,
ClinVar is not a complete truth set. Many variants remain variants of uncertain
significance, and submitted interpretations can be sparse, conflicting, or
dependent on the reference sequence used. For MAPT/Tau-F, strict
reference-residue matching leaves only a small binary benchmark.

Protein language models provide a complementary approach. These models are
trained on large collections of natural protein sequences and learn statistical
patterns related to evolutionary and biochemical constraints
\citep{rives2021biological}. ESM-1v showed that pretrained protein language
models can be used in a zero-shot setting to predict mutation effects without
task-specific experimental training data \citep{meier2021language}. General
missense predictors such as AlphaMissense provide another useful point of
comparison, but they also require coordinate quality control in isoform-rich
genes such as MAPT \citep{cheng2023alphamissense}.

Here, we present a reproducible zero-shot missense atlas for 441-aa Tau-F. We
enumerated all 8379 possible single amino-acid substitutions, scored them with an
ESM-1v ensemble, added Tau-specific mechanistic annotations, and integrated
ClinVar and AlphaMissense through strict reference-residue quality control.

\section{Methods}

\subsection{Study design}

This study builds a label-free missense variant atlas for MAPT/Tau. We first
listed every possible single amino-acid change in Tau-F, then used pretrained
protein language models to score those changes without using clinical labels for
training or tuning. Clinical databases were added only after scoring, as
evaluation and interpretation layers.

\subsection{Tau reference sequence and variant enumeration}

All analyses used the 441-amino-acid adult CNS Tau-F / hTau40 / 2N4R isoform as
the reference sequence. For each of the 441 positions, all 19 possible
non-reference amino-acid substitutions were generated. This produced 8379
missense variants. Each variant was assigned a compact identifier such as P301L,
meaning the reference amino acid proline at position 301 is changed to leucine.

\subsection{Tau-specific annotation}

Each variant was annotated with interpretable Tau biology features, including Tau
region, aggregation motif proximity, post-translational modification proximity,
known pathogenic hotspot proximity, charge-changing substitutions, and changes
involving special residues such as proline, glycine, or cysteine. These
annotations were not used to train ESM. They were used for interpretation,
prioritization, and a transparent heuristic baseline.

\subsection{Zero-shot ESM scoring}

Variants were scored with pretrained ESM protein language models using masked
marginal scoring. For each Tau-F position, the model estimated the log
probability of the reference amino acid and the mutant amino acid in the same
sequence context. The raw log-likelihood ratio was
\[
  \mathrm{esm\_llr} = \log p(\mathrm{mutant}) - \log p(\mathrm{wildtype}).
\]
For easier prioritization, we used
\[
  \mathrm{pathogenic\_score} = -\mathrm{esm\_llr}.
\]
Higher scores therefore indicate stronger zero-shot sequence incompatibility. The
manuscript-scale run used five ESM-1v checkpoints and averaged their scores.

\subsection{ClinVar and AlphaMissense quality control}

ClinVar and AlphaMissense annotations were imported after zero-shot scoring. Rows
were accepted only if the reported reference amino acid matched the 441-aa Tau-F
reference sequence at the mapped position. Rows outside the 1--441 range or with
reference-residue mismatches were rejected and written to quality-control tables.
This strict matching prevents accidental mixing of MAPT isoform coordinate
systems.

\subsection{Heuristic baseline and calibration}

A transparent Tau heuristic baseline was created to test whether ESM scores only
reflected obvious Tau-domain features. ESM-1v and the heuristic were compared
across all 8379 variants using a top-decile rule. AlphaMissense was retained as a
supplementary coverage-limited external reference rather than a primary full-atlas
comparator. Established MAPT controls were ranked against the complete ESM atlas.
\section{Results}

\subsection{A complete Tau-F missense atlas}

We generated a complete single amino-acid substitution atlas for the 441-residue
2N4R Tau-F / hTau40 reference sequence. The atlas contains 8379 missense
variants, corresponding to 19 possible substitutions at each reference position.
Clinical annotations were added only after the full variant list and zero-shot
model scores had been generated.

\subsection{Strict clinical-label coordinate QC limits direct benchmarking}

ClinVar import retained 33 annotated variants and rejected 1119 MAPT rows that
did not cleanly map to the Tau-F atlas. Accepted labels were sparse: one
pathogenic or likely pathogenic variant, two benign or likely benign variants,
26 variants of uncertain significance, and four variants with conflicting
classifications. Therefore, the direct P/LP versus B/LB benchmark contains only
three examples. AUROC and AUPRC values from this benchmark should be interpreted
as pipeline checks rather than clinical performance estimates.

\subsection{ESM-1v ensemble scores highlight Tau repeat regions}

At the domain level, the highest average ESM-1v ensemble scores were observed in
the microtubule repeat regions, while the N-terminal projection domain had the
lowest average score (Table~\ref{tab:domain}).

\begin{table}[H]
\centering
\caption{Domain-level ESM-1v ensemble summary.}
\label{tab:domain}
\begin{tabular}{lrr}
\toprule
Region & Mean score & Top 10\% fraction \\
\midrule
microtubule repeat R4 & 13.22 & 0.262 \\
microtubule repeat R3 & 13.21 & 0.255 \\
microtubule repeat R2 / exon 10 & 12.96 & 0.248 \\
microtubule repeat R1 & 12.64 & 0.211 \\
C-terminal tail & 12.03 & 0.114 \\
proline-rich region & 9.83 & 0.057 \\
N-terminal projection & 3.08 & 0.000 \\
\bottomrule
\end{tabular}
\end{table}

\subsection{Model concordance and calibration}

Using the top-decile rule across the complete atlas, 209 variants were high-priority
for both ESM-1v and the Tau heuristic, while 629 were ESM-only and 629 were
heuristic-only. We separately ranked five established MAPT missense controls:
G272V, P301L, V337M, R406W, and N279K. None entered the ESM-1v top 10%. G272V
ranked 2097/8379 (top 25.03%), P301L 3161/8379 (top 37.73%), N279K 3076/8379
(top 36.71%), R406W 3184/8379 (top 38.00%), and V337M 3870/8379 (top 46.19%).
This calibration result limits clinical interpretation of the score and supports its
use as a prioritization layer rather than a calibrated pathogenicity probability.
\subsection{ClinVar VUS prioritization}

Because most strictly mapped ClinVar variants were VUS, prioritization is more
appropriate than binary classification for the current clinical layer. The top
ESM-1v-ranked ClinVar VUS candidates included G440E, G333A, K290Q, T377A, and
D34Y (Table~\ref{tab:vus}). These variants should be treated as candidates for
follow-up review, not as clinical reclassifications.

\begin{table}[H]
\centering
\caption{Top ESM-1v-ranked ClinVar VUS candidates.}
\label{tab:vus}
\begin{tabular}{rlrl}
\toprule
Rank & Variant & Score & Region \\
\midrule
1 & G440E & 10.37 & C-terminal tail \\
2 & G333A & 8.73 & microtubule repeat R3 \\
3 & K290Q & 8.48 & microtubule repeat R2 / exon 10 \\
4 & T377A & 7.02 & C-terminal tail \\
5 & D34Y & 4.12 & N-terminal projection \\
\bottomrule
\end{tabular}
\end{table}

\subsection{AlphaMissense coverage}

AlphaMissense scores were available for 877 of 8379 Tau-F missense variants,
covering 149 positions. Among scored variants, 798 were AlphaMissense likely
benign, 67 ambiguous, and 12 likely pathogenic. The importer rejected 3932 MAPT
AlphaMissense rows after filtering for P10636: 2048 were outside the 441-aa
Tau-F reference range, and 1884 had a reference amino-acid mismatch.

\section{Discussion}

This study produced a complete zero-shot missense atlas for the 441-aa Tau-F /
hTau40 isoform. The atlas combines ESM-1v ensemble scores with Tau-specific
mechanistic annotations, positive-control calibration, ClinVar coordinate QC,
supplementary AlphaMissense QC, and ranked VUS prioritization. The strongest claim
is not clinical diagnosis. The strongest claim is that the atlas gives researchers
a transparent starting point for MAPT/Tau variant review, hypothesis generation,
and follow-up prioritization.

The ESM-1v ensemble assigned higher average scores to the microtubule repeat
regions than to the N-terminal projection domain. This pattern is biologically
plausible because the repeat regions are central to Tau microtubule binding and
aggregation-related biology. However, the five-control calibration showed that
known MAPT missense controls were not concentrated in the ESM-1v top decile. The
regional signal should therefore be interpreted as a prioritization pattern, not
proof that every high-scoring substitution is pathogenic.

A separate calibration analysis exposed an important limitation. Five established
MAPT missense controls (G272V, P301L, V337M, R406W, and N279K) were all present in
the atlas, but none was placed in the ESM-1v top decile. G272V was the highest-ranked
control at 2097/8379, corresponding to the top 25.03% of scores; the remaining
controls fell between the top 36.71% and top 46.19%. The present ESM-1v score should
therefore not be described as a calibrated MAPT pathogenicity probability. It is
better interpreted as a sequence-compatibility prioritization signal whose usefulness
depends on independent Tau biology and variant-level evidence.

A major practical lesson is that coordinate QC is not a minor technical detail.
Strict ClinVar and AlphaMissense matching showed that many external MAPT rows
cannot be safely mapped to the 441-aa Tau-F reference. AlphaMissense covered only
877 of 8379 atlas variants, so it is retained as a supplementary reference rather
than a primary full-atlas comparator. This does not mean that ClinVar or
AlphaMissense are low-quality resources; it means that MAPT isoform biology makes
naive merging risky.

The strict binary ClinVar benchmark contains only three P/LP versus B/LB examples.
Therefore, AUROC and AUPRC values should be treated as smoke checks rather than
clinical-performance estimates. Current VUS rankings are also hypothesis-generating,
not clinical interpretations. G440E is a concrete example: although it lies in the
C-terminal tail rather than a canonical microtubule-binding repeat, its high score
could reflect sequence constraints related to Tau conformational dynamics, membrane
association, or phosphorylation accessibility. This hypothesis should be tested in
cellular or biochemical assays rather than treated as a conclusion.

Several limitations remain. ESM-1v does not directly model splicing, isoform
expression, post-translational modification state, neuronal cell context,
aggregation kinetics, or patient-level genetic background. The negative positive-
control calibration indicates that a generic protein language model may not rank
every known Tau disease mutation highly. AlphaMissense coverage is incomplete, and
all current VUS candidates require independent evidence and experimental validation.

Future work should add population-frequency filtering, conservation and disorder
features, and functional assays. The most useful next computational analysis is a
mechanistic regression asking whether ESM scores are associated with conservation,
hydrophobicity, disorder, charge change, or aggregation-related features.
\section{Figures}

\begin{figure}[H]
\centering
\includegraphics[width=0.95\linewidth]{figure_workflow_schematic.png}
\caption{MAPT/Tau zero-shot atlas workflow. Clinical labels and external
predictors are added after label-free atlas scoring, with strict coordinate QC.}
\end{figure}

\begin{figure}[H]
\centering
\includegraphics[width=0.95\linewidth]{missense_heatmap.png}
\caption{ESM-1v missense score heatmap across Tau-F. Each position corresponds to
one amino-acid site and each row to a mutant amino acid.}
\end{figure}

\begin{figure}[H]
\centering
\includegraphics[width=0.95\linewidth]{score_by_position.png}
\caption{ESM-1v ensemble scores by Tau-F position.}
\end{figure}

\begin{figure}[H]
\centering
\includegraphics[width=0.85\linewidth]{figure_domain_summary.png}
\caption{Domain-level ESM-1v score patterns across Tau-F regions.}
\end{figure}

\begin{figure}[H]
\centering
\includegraphics[width=0.85\linewidth]{figure_model_comparison.png}
\caption{Strict ClinVar smoke benchmark comparison. The benchmark is too small
for clinical performance claims.}
\end{figure}

\begin{figure}[H]
\centering
\includegraphics[width=0.85\linewidth]{figure_esm1v_vs_heuristic.png}
\caption{ESM-1v ensemble scores compared with the transparent Tau heuristic
baseline across all 8379 variants.}
\end{figure}

\begin{figure}[H]
\centering
\includegraphics[width=0.95\linewidth]{figure_vus_lollipop.png}
\caption{Top ClinVar VUS candidates across Tau-F, ranked by ESM-1v ensemble
pathogenic score.}
\end{figure}

\begin{figure}[H]
\centering
\includegraphics[width=0.82\linewidth]{figure_positive_control_calibration.png}
\caption{Calibration against five established MAPT missense controls. None of the
controls entered the ESM-1v top 10%, so the score is used as prioritization evidence
rather than a clinically calibrated pathogenicity probability.}
\label{fig:positive_controls}
\end{figure}
\section{Data and code availability}

All source code for the MAPT/Tau zero-shot atlas workflow is available at
\url{https://github.com/xiaopangF/Tau-zero-shot-MAPT}. Large generated files are
not committed directly to GitHub. They can be regenerated from public ClinVar,
AlphaMissense, and ESM resources using the repository workflow. Before journal
submission, generated result tables and figures should be archived with a
permanent DOI through a repository such as Zenodo or Figshare.

\section{Supplementary tables}

Supplementary tables include the full Tau-F missense atlas, ESM-1v ensemble
scores, ClinVar coordinate QC tables, AlphaMissense coordinate QC tables, the top
VUS priority list, and model concordance/discordance tables.

\bibliographystyle{unsrtnat}
\bibliography{../references}

\end{document}